\clearpage
\section*{Problem 3}
\addcontentsline{toc}{section}{Problem 3}
\begingroup
\hypersetup{colorlinks=true,
            linkcolor=blue!50!black,
            citecolor=blue!50!black,
            urlcolor=blue!50!black}
\newcommand{\qthreeR}{\mathbb{R}}
\newcommand{\qthreeN}{\mathbb{N}}
\newcommand{\qthreeZ}{\mathbb{Z}}
\newcommand{\qthreeSnl}{S_n(\lambda)}
\newcommand{\qthreewt}{\operatorname{wt}}
\newtheorem{qthreetheorem}{Theorem}
\newtheorem{qthreelemma}{Lemma}
\newtheorem{qthreeproposition}{Proposition}
\newtheorem{qthreecorollary}{Corollary}
\theoremstyle{definition}
\newtheorem{qthreedefinition}{Definition}
\newtheorem{qthreeremark}{Remark}
\title{Solution to Problem 3: A Probabilistic Interpretation\\
for Interpolation Macdonald Polynomials}
\author{}
\date{}
\maketitle


\section*{Statement and Answer}

Let $\lambda=(\lambda_1>\cdots>\lambda_n\ge0)$ be a partition with distinct
parts that is moreover \emph{restricted}: it has a unique part of size $0$ and no
part of size $1$. Let $\qthreeSnl$ be the orbit of $\lambda$ under permutation of
parts (equivalently, the compositions that rearrange $\lambda$). Let
$F^*_\mu(x_1,\dots,x_n;q,t)$ and $P^*_\lambda(x_1,\dots,x_n;q,t)$ be the
interpolation ASEP polynomial and interpolation Macdonald polynomial.

\medskip
\noindent\textbf{Theorem.} \emph{There is a nontrivial Markov chain on $\qthreeSnl$ ---
the \emph{interpolation $t$-Push TASEP} --- whose stationary distribution is}
\[
\pi^*_\lambda(\mu)\;=\;\frac{F^*_\mu(x_1,\dots,x_n;\,q=1,\,t)}{P^*_\lambda(x_1,\dots,x_n;\,q=1,\,t)},
\qquad \mu\in\qthreeSnl,
\]
\emph{and whose transition probabilities are not described using the
polynomials $F^*_\mu$.}

\medskip
\noindent The answer is \textbf{YES}. We caution at the outset that the chain is
\emph{not} the ordinary multispecies ASEP: that chain has the ordinary (not
interpolation) Macdonald polynomials as its stationary weights (Ayyer--Martin--Williams),
and substituting it here would solve a different, already-known problem. The
correct object is a $q=1$ \emph{push}-TASEP with an interpolation correction
encoded by a ringing ``bell''.

\section*{Probabilities and notation}

Throughout we work at $q=1$ and write $P^*_\lambda=P^*_\lambda(x;1,t)$, etc.
For $1\le k\le n$ put
\[
p_k:=\frac{t^{-n+1}(1-t)}{x_k-t^{-n+2}},\qquad
q_k:=\frac{(1-t)\,x_k}{x_k-t^{-n+2}} .
\]
For $0<t<1$ and $x_i>t^{-n+1}$ these lie in $[0,1]$ and serve as
skipping/displacement probabilities. We use $e^*_k(x;t)$ for the interpolation
elementary symmetric polynomials and the factorization (valid at $q=1$
\cite{Dolega,BDW25a})
\begin{equation}\label{eq:factor}
P^*_\lambda(x;1,t)=\prod_{1\le i\le\lambda_1}e^*_{\lambda'_i}(x;t),
\end{equation}
$\lambda'$ the conjugate partition.

We record two elementary identities used repeatedly below; they are the
algebraic content of the correspondence in Proposition~\ref{prop:step2}. Since
$t^{-n+2}-t^{-n+1}(1-t)=t^{-n+1}\bigl(t-(1-t)\bigr)$ is \emph{not} what we want,
we compute directly:
\begin{align}
\bigl(x_k-t^{-n+2}\bigr)-t^{-n+1}(1-t)
   &=x_k-t^{-n+2}-t^{-n+1}+t^{-n+2}=x_k-t^{-n+1},
   \label{eq:idp}\\
\bigl(x_k-t^{-n+2}\bigr)-(1-t)\,x_k
   &=t\,x_k-t^{-n+2}=t\bigl(x_k-t^{-n+1}\bigr).
   \label{eq:idq}
\end{align}
Consequently, dividing \eqref{eq:idp} and \eqref{eq:idq} by $x_k-t^{-n+2}$,
\begin{equation}\label{eq:skipprobs}
1-p_k=\frac{x_k-t^{-n+1}}{x_k-t^{-n+2}},\qquad
1-q_k=\frac{t\,\bigl(x_k-t^{-n+1}\bigr)}{x_k-t^{-n+2}} .
\end{equation}

\section*{The interpolation $t$-Push TASEP}

\begin{qthreedefinition}\label{def:chain}
Fix $\lambda$ with $\lambda_n=0$. States of the chain are configurations of
particles on a ring of $n$ sites labelled by the parts of $\lambda$: state $\eta$
places the particle labelled $\eta_j$ at site $j$ (label $0$ is a vacancy). One
step from $\eta$ proceeds in three stages.

\smallskip
\noindent\textbf{(Step 0) The bell.} The bell at site $j$ rings with probability
\[
P_j=\frac{\prod_{k<j}\bigl(x_k-t^{-n+2}\bigr)\,\prod_{k>j}\bigl(x_k-t^{-n+1}\bigr)}
{e^*_{n-1}(x;t)},
\quad
e^*_{n-1}(x;t)=\sum_{j=1}^n\prod_{k<j}\bigl(x_k-t^{-n+2}\bigr)\prod_{k>j}\bigl(x_k-t^{-n+1}\bigr).
\]

\smallskip
\noindent\textbf{(Step 1) Push-TASEP move.} The particle at site $j$, say with
label $a$, is activated and travels clockwise by the $t$-Push TASEP rule: if
there are $m$ weaker particles in the system (labels $<a$, vacancies counting as
label $0$), the active particle moves to the $k$-th such weaker particle with
probability $t^{k-1}/[m]_t$. If that site held a positive label, that particle
becomes active and repeats; the cascade ends when the active particle reaches a
vacancy. Afterwards site $j$ is vacant; we treat this vacancy as a new particle
labelled $a:=0$.

\smallskip
\noindent\textbf{(Step 2) Vacancy sweep.} The label-$0$ particle goes to site $1$
and travels clockwise. Reaching a site $k$ ($1\le k\le j-1$) holding a label
$b\ge0$, it skips that site with probability $1-p_k$ if $b\ge a$ and $1-q_k$ if
$b<a$; otherwise it settles there, activating/displacing the resident, which then
continues clockwise toward site $j$ by the same rule. The active particle stops
once it displaces another or arrives at site $j$, where it settles.
\end{qthreedefinition}

We denote the transition probabilities of Steps 1 and 2 by $P^{(1)}_j$ and
$P^{(2)}_j$, so that for $\mu,\nu\in\qthreeSnl$,
\[
P(\mu,\nu)=\sum_{1\le j\le n}P_j\sum_{\rho\in\qthreeSnl:\rho_j=0}
P^{(1)}_j(\mu,\rho)\,P^{(2)}_j(\rho,\nu).
\]

\begin{qthreeremark}[nontriviality]
The rates $P_j,p_k,q_k$ are explicit rational functions of $x_1,\dots,x_n$ and
$t$ only. They do \emph{not} reference $F^*_\mu$, so the chain is nontrivial in
the required sense (unlike a Metropolis--Hastings chain built from the target
distribution).
\end{qthreeremark}

\begin{qthreetheorem}\label{thm:main}
In the interpolation $t$-Push TASEP with content $\lambda$, the stationary
probability of $\mu\in\qthreeSnl$ is $\pi^*_\lambda(\mu)=F^*_\mu(x;1,t)/P^*_\lambda(x;1,t)$.
\end{qthreetheorem}

\section*{Two-line queues and the proof}

We use classical two-line queues \cite{CMW22} and signed two-line queues
\cite{BDW25a}, specialized to $q=1$. Let $\mathcal Q^\eta_\kappa$ be the classical
two-line queues with top row $\eta$ and bottom row $\kappa$, with pair-weight
generating function $a^\eta_\kappa=\sum_{Q\in\mathcal Q^\eta_\kappa}\qthreewt_{\mathrm{pair}}(Q)$.
Let $\mathcal G^\alpha_\mu$ be the signed two-line queues with top row $\alpha\in\qthreeZ^n$
and bottom row $\mu$, with $b^\alpha_\mu=\sum_{Q\in\mathcal G^\alpha_\mu}\qthreewt_{\mathrm{pair}}(Q)$,
and full weight $\qthreewt(Q)=\qthreewt_{\mathrm{pair}}(Q)\qthreewt_{\mathrm{ball}}(Q)$ where
\begin{equation}\label{eq:wta}
\qthreewt_\alpha\,b^\alpha_\mu=\sum_{Q\in\mathcal G^\alpha_\mu}\qthreewt(Q),
\quad \qthreewt_\alpha:=\prod_{k:\alpha_k>0}x_k\prod_{k:\alpha_k<0}\bigl(-t^{-(n-1)}\bigr).
\end{equation}

\begin{qthreedefinition}\label{def:unsigned}
For a signed two-line queue $Q\in\mathcal G^\alpha_\mu$ let $\overline Q$ be its
unsigned version (forget signs on the top row); its top composition is
$\|\alpha\|=(|\alpha_1|,\dots,|\alpha_n|)$. Let $\mathcal G^\kappa_\mu$ be the set of
unsigned objects so obtained, with $\|\alpha\|=\kappa$. For $\overline Q\in\mathcal G^\kappa_\mu$
assign a nontrivial pairing $p$ the weight $(1-t)t^{\mathrm{skip}(p)}$, where
$\mathrm{skip}(p)$ is the number of weaker balls (labels strictly smaller than the
moving label, vacancies included) over which $p$ jumps, and a top-row
ball $B$ in column $k$ labelled $a>0$ with the ball below labelled $b$ the weight
\[
\qthreewt(B)=\begin{cases}x_k-t^{-(n-1)} & b=a,\\ x_k & b>a,\\ t^{-(n-1)} & b<a.\end{cases}
\]
\end{qthreedefinition}

\begin{qthreelemma}\label{lem:signforget}
For $\lambda$ with distinct parts and $\kappa,\mu\in\qthreeSnl$, and $\overline Q\in\mathcal G^\kappa_\mu$,
$\qthreewt(\overline Q)=\sum_Q\qthreewt(Q)$, the sum over signed $Q$ forgetting to $\overline Q$.
\end{qthreelemma}

\begin{proof}
Enumerate the admissible sign choices for each top-row ball $B$ labelled $a>0$:
a $-$ sign is forced when the ball below is a vacancy or labelled $b<a$; a $+$
sign is forced when $b>a$; both are allowed when $b=a$. One checks this matches
\eqref{eq:wta}: a $-$ sign multiplies the ball weight by $-1$ in passing from
$Q$ to $\overline Q$, but the weight of the pairing attached to $B$ is also
multiplied by $-1$, so signs cancel consistently.
\end{proof}

For $\kappa\in\qthreeSnl$ set $c^\kappa_\mu:=\sum_{\alpha:\|\alpha\|=\kappa}\qthreewt_\alpha\,b^\alpha_\mu$.
Combining \eqref{eq:wta} and Lemma~\ref{lem:signforget}:

\begin{qthreelemma}\label{lem:ckm}
$c^\kappa_\mu=\sum_{\overline Q\in\mathcal G^\kappa_\mu}\qthreewt(\overline Q)$. Since $\lambda$ has distinct
parts, $\mathcal G^\kappa_\mu$ is empty or a single element.
\end{qthreelemma}

\subsection*{The closing recursion (where the hypotheses on $\lambda$ enter)}

Fix a weakly order-preserving $\phi:\qthreeN\to\qthreeN$ and partitions $\lambda,\kappa$ with
$\phi(\lambda)=\kappa$. For $\eta\in S_n(\kappa)$ set
$G^*_\eta(x;t):=\sum_{\rho\in\qthreeSnl:\,\phi(\rho)=\eta}F^*_\rho(x;1,t)$. The following
is an interpolation analogue of a result of Ayyer--Martin--Williams
\cite{AMW25}, provable in the same way using \cite{AS19}.

\begin{qthreetheorem}\label{thm:G}
For $\lambda,\kappa$ as above and all $\eta\in S_n(\kappa)$, at $q=1$,
\[
\frac{G^*_\eta(x;t)}{P^*_\lambda(x;1,t)}=\frac{F^*_\eta(x;1,t)}{P^*_\kappa(x;1,t)} .
\]
\end{qthreetheorem}

\begin{qthreecorollary}\label{cor:rec}
Let $\rho$ be a composition with $\rho_i\ne1$ for all $i$, with $k$ nonzero
parts, and set $\eta=\rho^-$ (where $\rho^-_i=\max(\rho_i-1,0)$). Then at $q=1$,
\[
F^*_\rho(x;1,t)=F^*_\eta(x;1,t)\cdot e^*_k(x;t).
\]
\end{qthreecorollary}

\begin{proof}
Let $\lambda,\kappa$ reorder $\rho,\eta$. Take $\phi:i\mapsto\max(i-1,0)$, so
$\phi(\rho)=\eta$. \emph{Because $\lambda$ has no part of size $1$}, $\phi$ is a
bijection from $\{0,2,3,\dots\}$ onto $\{0,1,2,\dots\}$, so $\rho$ is the unique
composition in $\qthreeSnl$ with $\phi(\rho)=\eta$; hence $G^*_\eta=F^*_\rho$. By
Theorem~\ref{thm:G},
$F^*_\rho/P^*_\lambda=F^*_\eta/P^*_\kappa$. Finally, by \eqref{eq:factor} and the
fact that $\kappa$ is obtained from $\lambda$ by deleting the largest column (of
size $k$), $P^*_\lambda/P^*_\kappa=e^*_k(x;t)$, giving the claim.
\end{proof}

This is exactly where the restriction on $\lambda$ is essential: \emph{no part of
size $1$} makes $\phi$ bijective so the recursion closes, and the \emph{unique
part of size $0$} guarantees $\eta$ has a unique $0$-part so the column-deletion
is unambiguous. The hypotheses are structural, not (as one might guess) a device
to avoid a $0/0$ normalization.

\subsection*{From two-line queues to the dynamics}

\begin{qthreeproposition}\label{prop:step2}
For $\rho,\nu\in\qthreeSnl$ with $\rho_j=0$,
\[
P^{(2)}_j(\rho,\nu)=\frac{c^\rho_\nu}{\prod_{k<j}(x_k-t^{-n+2})\prod_{k>j}(x_k-t^{-n+1})},
\quad\text{equivalently}\quad
P_j\,P^{(2)}_j(\rho,\nu)=\frac{c^\rho_\nu}{e^*_{n-1}} .
\]
\end{qthreeproposition}

\begin{proof}
Write $D:=\prod_{k<j}(x_k-t^{-n+2})\prod_{k>j}(x_k-t^{-n+1})$, so that $P_jD=
c^\rho_\nu/e^*_{n-1}$ is the second assertion once the first is proved. Since
$\rho$ rearranges $\lambda$, which has distinct parts, $\mathcal G^\rho_\nu$ is empty
or a single element $\overline Q$ (Lemma~\ref{lem:ckm}), and $c^\rho_\nu=\qthreewt(\overline Q)$
in the latter case, $c^\rho_\nu=0$ in the former. We prove
\begin{equation}\label{eq:goal}
\qthreewt(\overline Q)=D\cdot P^{(2)}_j(\rho,\nu)\qquad(\overline Q\in\mathcal G^\rho_\nu),
\qquad\text{and}\qquad P^{(2)}_j(\rho,\nu)=0\ \text{when}\ \mathcal G^\rho_\nu=\varnothing,
\end{equation}
which is \eqref{eq:goal} in both cases.

\medskip
\noindent\emph{Step A: the dynamics is a tree of skip/settle decisions.}
In Step 2 a single active label $a$ is carried clockwise across the sites
$k=1,2,\dots,j-1$ (sites $k>j$ are never visited; site $j$ is the terminal site).
Initially $a=0$. On reaching a site $k<j$ holding a resident $b$ the chain makes
one of two moves:
\begin{itemize}
\item[(\emph{skip})] the resident $b$ stays at site $k$ and the \emph{same} label
$a$ continues to site $k+1$; the probability is $1-p_k$ if $b\ge a$ and $1-q_k$
if $b<a$;
\item[(\emph{settle})] the label $a$ stops at site $k$ (so the new content of
site $k$ is $a$), the resident $b$ becomes the new active label and continues to
site $k+1$; the probability is $p_k$ if $b\ge a$ and $q_k$ if $b<a$.
\end{itemize}
At site $j$ (whose original content is $\rho_j=0$) the current active label
settles, ending the step. Thus every realisation of Step 2 is a path in a binary
tree whose nodes are the visited sites $1,\dots,j-1$; the leaf at site $j$
determines the output configuration $\nu$, and
$P^{(2)}_j(\rho,\nu)$ is the sum over all root-to-leaf paths producing $\nu$ of
the product of the chosen branch probabilities.

\medskip
\noindent\emph{Step B: the weight emitted on each branch, divided by the matching
factor of $D$, is the branch probability.}
We assign to each branch a weight that is a single ball weight (and, for
settles, a factor $1-t$), and show that the emitted weight divided by the
corresponding factor of $D$ is exactly the branch probability listed in Step~A.
There are four branches, governed by the resident relation $b\ge a$ or $b<a$.
The matching factor of $D$ at a visited site $k<j$ is $x_k-t^{-n+2}$.

\smallskip
\noindent(1) \emph{skip, $b\ge a$.} The resident $b$ remains at site $k$, so in
the two-line queue the column-$k$ top ball (content of site $k$ \emph{before} the
move) and bottom ball (content \emph{after}) carry the same label $b$, i.e.\ the
case ``$b=a$'' of Definition~\ref{def:unsigned}; its ball weight is
$x_k-t^{-(n-1)}=x_k-t^{-n+1}$. Dividing by $x_k-t^{-n+2}$ gives, by the first
identity in \eqref{eq:skipprobs}, exactly $1-p_k$. We therefore emit the ball
weight $x_k-t^{-n+1}$ on this branch.

\smallskip
\noindent(2) \emph{skip, $b<a$.} Here the moving label $a$ passes over a strictly
weaker resident $b$. The resident again stays put (same-label column, ball weight
$x_k-t^{-n+1}$ as in (1)), but the moving label has jumped one weaker ball,
contributing one factor $t$ to the eventual pairing weight $(1-t)t^{\mathrm{skip}}$
of the pairing carrying $a$. Hence the emitted weight is $t\,(x_k-t^{-n+1})$, and
dividing by $x_k-t^{-n+2}$ gives, by the second identity in \eqref{eq:skipprobs},
exactly $1-q_k$.

\smallskip
\noindent(3) \emph{settle, $b\ge a$.} The label $a$ stops at site $k$ and the
resident $b\ge a$ is displaced. In the queue this opens a nontrivial pairing
that begins at column $k$ (the pairing that will carry $b$ onward); the new
column-$k$ ball relation is the displaced label below a smaller settled label,
which is the case ``$b<a$'' of Definition~\ref{def:unsigned}, of ball weight
$t^{-(n-1)}=t^{-n+1}$. The pairing just opened contributes its base factor
$1-t$ (its $t^{\mathrm{skip}}$ accumulates on later skip-branches of type (2)).
Thus the emitted weight is $t^{-n+1}(1-t)$, and dividing by $x_k-t^{-n+2}$ gives
exactly $p_k=t^{-n+1}(1-t)/(x_k-t^{-n+2})$.

\smallskip
\noindent(4) \emph{settle, $b<a$.} The label $a$ stops at site $k$ and the
strictly weaker resident $b<a$ is displaced; the new column-$k$ ball relation is
the displaced label below a larger settled label, the case ``$b>a$'' of
Definition~\ref{def:unsigned}, of ball weight $x_k$. As in (3) a nontrivial
pairing opens, contributing $1-t$. The emitted weight is $x_k(1-t)$, and dividing
by $x_k-t^{-n+2}$ gives exactly $q_k=(1-t)x_k/(x_k-t^{-n+2})$.

\smallskip
In all four cases the emitted weight at a visited site $k<j$, divided by the
factor $x_k-t^{-n+2}$ of $D$, equals the branch probability.

\medskip
\noindent\emph{Step C: the terminal site and the sites $k>j$.}
At the terminal site $j$ the active label settles; the original content there is
$\rho_j=0$, so column $j$ has a top vacancy and emits no ball weight (and is not a
factor of $D$). The sites $k>j$ are never visited by the sweep, so their content
is unchanged: $\nu_k=\rho_k$. In the queue, column $k>j$ then has equal top and
bottom labels, the case ``$b=a$'', of ball weight $x_k-t^{-(n-1)}=x_k-t^{-n+1}$,
which is precisely the matching factor of $D$ for $k>j$. Hence the column-$k$
contribution divided by its $D$-factor equals $1$, consistent with the statement
that columns $k>j$ never move.

\medskip
\noindent\emph{Step D: telescoping a settle-move into one pairing.}
A nontrivial pairing in $\overline Q$ corresponds to a maximal run in which a label
is carried forward: it is opened by a settle-branch (type (3) or (4)), then
crosses some number $s\ge0$ of weaker residents via skip-branches of type (2)
(each contributing one factor $t$), and finally settles again or reaches site
$j$. The factors accumulated on this run are $(1-t)$ from the opening settle and
$t$ from each of the $s$ type-(2) skips it traverses, giving total pairing weight
$(1-t)t^{s}=(1-t)t^{\mathrm{skip}(p)}$, in agreement with
Definition~\ref{def:unsigned}; here $\mathrm{skip}(p)=s$ is the number of weaker
balls the pairing jumps over. Each type-(1) skip and each unvisited site $k>j$
contributes its own same-label ball weight $x_k-t^{-n+1}$ as in (1) and Step~C.
Multiplying the contributions of all columns therefore yields exactly
$\qthreewt(\overline Q)$ on the one hand (product of the ball weights of all $n$
columns times $(1-t)t^{\mathrm{skip}}$ over all nontrivial pairings) and, on the
other hand, the product of all branch probabilities of the realising path times
$D$ (each visited factor $x_k-t^{-n+2}$ for $k<j$ supplied by Step~B and each
factor $x_k-t^{-n+1}$ for $k>j$ supplied by Step~C). That is, for a single
root-to-leaf path $\sigma$ realising $\nu$,
\[
\qthreewt_\sigma(\overline Q)=D\cdot(\text{product of branch probabilities along }\sigma).
\]

\medskip
\noindent\emph{Step E: uniqueness of the realising path and the empty case.}
A clockwise sweep moves each label only from its site toward site $j$ (a settle
fixes the active label at the current site and hands the resident forward) or
leaves it fixed (a skip). Consequently the output $\nu$ determines the path
$\sigma$ uniquely: reading $\nu$ from site $1$ to site $j$, the label $\nu_k$ at a
visited site $k<j$ was either the resident that stayed (skip at $k$) or the active
label that settled (settle at $k$), and which of the two occurred is forced by
whether $\nu_k=\rho_k$ or not, since $\lambda$ has distinct parts so the resident
and the active label are never equal except in a skip with $\nu_k=\rho_k=b$.
Thus there is \emph{exactly one} root-to-leaf path $\sigma$ producing a given
realisable $\nu$. Moreover the column ball-cases read off from $\sigma$ in
Steps~B--C are exactly the column ball-cases of the unique
$\overline Q\in\mathcal G^\rho_\nu$ (Lemma~\ref{lem:ckm}), and its nontrivial pairings
are exactly the settle-runs of $\sigma$ described in Step~D; hence
$\qthreewt_\sigma(\overline Q)=\qthreewt(\overline Q)=c^\rho_\nu$. Dividing the identity
$\qthreewt_\sigma(\overline Q)=D\cdot(\text{branch product along }\sigma)$ of Step~D
by $D$,
\[
P^{(2)}_j(\rho,\nu)=(\text{branch product along }\sigma)
=\frac{\qthreewt(\overline Q)}{D}=\frac{c^\rho_\nu}{D},
\]
which is \eqref{eq:goal}. Finally, if no two-line queue has top $\rho$ and bottom
$\nu$ (so $\mathcal G^\rho_\nu=\varnothing$ and $c^\rho_\nu=0$), then no sequence of
skip/settle moves transforms $\rho$ into $\nu$, because the realisable images of
$\rho$ under the sweep are precisely the bottom rows $\nu$ admitting such a queue
(each settle-run is a forward transport of a label, which is exactly a nontrivial
pairing, and the fixed labels are exactly the trivial same-label columns); hence
$P^{(2)}_j(\rho,\nu)=0=c^\rho_\nu/D$. This proves \eqref{eq:goal} in all cases,
and with it the proposition.

Finally, summing \eqref{eq:goal} over all $\nu$ (with $D$ fixed) and using that
the branch probabilities at each visited site sum to $1$
($p_k+(1-p_k)=q_k+(1-q_k)=1$) shows $\sum_\nu P^{(2)}_j(\rho,\nu)=1$, so
$P^{(2)}_j(\rho,\cdot)$ is genuinely a probability distribution; combined with
$0<t<1$ and $x_i>t^{-n+1}$ (whence $p_k,q_k\in[0,1]$ by \eqref{eq:skipprobs} and
positivity of the numerators), all branch probabilities lie in $[0,1]$.
\end{proof}

\begin{qthreeproposition}\label{prop:transition}
If $\lambda$ is restricted and $\mu,\nu\in\qthreeSnl$, then
$\displaystyle P(\mu,\nu)=\sum_{\rho\in\qthreeSnl}\frac{a^\mu_\rho\,c^\rho_\nu}{e^*_{n-1}}$.
\end{qthreeproposition}

\begin{proof}
Combine the encoding of Step 1 by classical two-line queues
\cite[Lem.~5.4]{AMW25} (giving $P^{(1)}_j(\mu,\rho)$ via $a^\mu_\rho$) with
Proposition~\ref{prop:step2}:
\[
P(\mu,\nu)=\sum_j P_j\!\!\sum_{\rho:\rho_j=0}\!P^{(1)}_j(\mu,\rho)P^{(2)}_j(\rho,\nu)
=\sum_j\sum_{\rho:\rho_j=0}\frac{a^\mu_\rho c^\rho_\nu}{e^*_{n-1}}
=\sum_{\rho\in\qthreeSnl}\frac{a^\mu_\rho c^\rho_\nu}{e^*_{n-1}}.
\]
\end{proof}

\subsection*{Proof of Theorem~\ref{thm:main}}

\begin{proof}
Fix $\nu\in\qthreeSnl$. By \cite[Thm.~1.15, Lem.~5.6]{BDW25a},
\[
F^*_\nu(x;1,t)=\sum_{\eta\in\qthreeN^n}F^{*\eta}_\nu(x;t)\,F^*_{\eta^-}(x;1,t),
\qquad
F^{*\eta}_\nu(x;t)=\sum_{\kappa\in\qthreeN^n}a^\eta_\kappa\,c^\kappa_\nu .
\]
By Corollary~\ref{cor:rec} (using that $\eta$ has a unique $0$-part),
$F^*_{\eta^-}(x;1,t)=F^*_\eta(x;1,t)/e^*_{n-1}(x;t)$. Substituting,
\[
F^*_\nu(x;1,t)=\sum_{\eta\in\qthreeN^n}F^*_\eta(x;1,t)\sum_{\kappa\in\qthreeN^n}
\frac{a^\eta_\kappa\,c^\kappa_\nu}{e^*_{n-1}(x;t)}
=\sum_{\eta\in\qthreeN^n}F^*_\eta(x;1,t)\,P(\eta,\nu),
\]
the last equality by Proposition~\ref{prop:transition}. Thus the vector
$\bigl(F^*_\mu(x;1,t)\bigr)_\mu$ is left-fixed by the transition matrix, i.e.\
proportional to the stationary distribution. Since
$P^*_\lambda=\sum_{\mu\in\qthreeSnl}F^*_\mu$, normalizing gives
$\pi^*_\lambda(\mu)=F^*_\mu(x;1,t)/P^*_\lambda(x;1,t)$.
\end{proof}

\section*{Remark on the failure mode to avoid}

Replacing this chain by the ordinary multispecies ASEP (matrix-product ansatz /
Zamolodchikov--Faddeev algebra) does \emph{not} solve the stated problem: that
construction yields the ordinary Macdonald polynomials as stationary weights
\cite{CMW22,AMW25}, not the \emph{interpolation} polynomials $F^*_\mu,P^*_\lambda$.
The interpolation correction is precisely what the bell $P_j$ and the
push-with-sweep dynamics supply, and the proof above runs through signed
two-line queues rather than a trace over an auxiliary algebra.

\begin{thebibliography}{9}
\bibitem{AS19} P.~Alexandersson, M.~Sawhney.
  \emph{Properties of non-symmetric Macdonald polynomials at $q=1$ and $q=0$}.
  Ann.\ Comb.\ \textbf{23} (2019), 219--239.
\bibitem{AMW25} A.~Ayyer, J.~Martin, L.~Williams.
  \emph{The inhomogeneous $t$-PushTASEP and Macdonald polynomials at $q=1$}.
  Ann.\ Inst.\ Henri Poincar\'e D, 2025.
\bibitem{BDW25a} H.~Ben Dali, L.~Williams.
  \emph{A combinatorial formula for interpolation Macdonald polynomials}.
  Preprint, arXiv:2510.02587, 2025.
\bibitem{BDW26} H.~Ben Dali, L.~Williams.
  \emph{A probabilistic interpretation for interpolation Macdonald polynomials}.
  Preprint, arXiv:2602.13492, 2026.
\bibitem{CMW22} S.~Corteel, O.~Mandelshtam, L.~Williams.
  \emph{From multiline queues to Macdonald polynomials via the exclusion process}.
  Amer.\ J.\ Math.\ \textbf{144} (2022), 395--436.
\bibitem{Dolega} M.~Dolega.
  \emph{Strong factorization property of Macdonald polynomials and higher-order
  Macdonald's positivity conjecture}.
  J.\ Algebraic Combin.\ \textbf{46} (2017), 135--163.
\end{thebibliography}

\endgroup
